\subsubsection{Граничные условия}

Для замыкания дискретной задачи необходимо задать граничные условия на давление~$p$ и (при использовании модели JFO) на степень заполнения~$\vartheta$. Тип граничных условий определяется физической постановкой (подразделы~2.1.2--2.1.4) и геометрией расчётной области. Ниже систематизированы три основных типа граничных условий для давления.

\textit{Дирихле (заданное давление).}
На соответствующем участке границы $\Gamma_D$ давление фиксируется:
\[
  p = p_0 \quad \text{на } \Gamma_D.
\]
Наиболее распространённый случай -- $p_0 = p_{\mathrm{cav}} = 0$ (избыточное давление) на осевых торцах подшипника, где смазочный слой свободно сообщается с окружающей средой.

\textit{Периодичность.}
По окружному (периодическому) направлению давление совпадает на противоположных границах области:
\[
  p(x_{1,s},\, x_2) = p(x_{1,t},\, x_2).
\]
Данное условие применяется в упорных и радиальных подшипниках, где расчётная область охватывает полный период по окружной координате.

\textit{Непроницаемость (Нейман).}
На участках с условием симметрии или непроницаемости нормальная производная давления обращается в нуль:
\[
  \frac{\partial p}{\partial n} = 0 \quad \text{на } \Gamma_N.
\]
Такое условие возникает, например, на линии симметрии при моделировании половины расчётной области.

В таблице~\ref{tab:bc_summary} сведены граничные условия для давления, соответствующие трём постановкам, сформулированным в подразделах~2.1.2--2.1.4.

\begin{table}[!ht]
\centering
\caption{Граничные условия для давления в различных постановках}
\label{tab:bc_summary}
\begin{tabular}{|l|l|l|}
\hline
\textbf{Постановка} & \textbf{Направление $x_1$} & \textbf{Направление $x_2$} \\
\hline
Упорная              & периодичность              & $p = 0$ на торцах          \\
\hline
Радиальная           & периодичность              & $p = 0$ на торцах          \\
\hline
Возвр.-поступательная & $p = 0$ на торцах          & периодичность или $p = 0$  \\
\hline
\end{tabular}
\end{table}

При использовании массо-сохраняющей модели кавитации JFO (раздел~2.3) дополнительно задаются граничные условия на степень заполнения~$\vartheta$. На границах подачи смазки полагается $\vartheta = 1$ (полное заполнение зазора). По периодическим направлениям степень заполнения периодична, аналогично давлению. На границах Дирихле ($p = p_{\mathrm{cav}}$) значение~$\vartheta$ не фиксируется, а определяется из уравнения сохранения массы (свободный выход смазки).

Во всех рассматриваемых постановках кавитационное давление принимается равным нулю в избыточных давлениях: $p_{\mathrm{cav}} = 0$.
